\documentclass[11pt]{article}
\usepackage[margin=1in]{geometry}
\usepackage{amsmath,amssymb}
\usepackage{booktabs}
\usepackage{graphicx}
\usepackage{hyperref}
\usepackage{xcolor}
\usepackage{natbib}
\usepackage{multirow}

\newcommand{\todo}[1]{\textcolor{red}{[TODO: #1]}}
\newcommand{\gain}{\mathrm{gain}}
\newcommand{\ncd}{\mathrm{NCD}}
\newcommand{\pm}{\mathrm{PM}}

\title{Mechanism-Matched Data Scoring:\\
Compression Predicts and Accelerates the Formation of Induction Heads}

\author{\todo{author, affiliation, email}}
\date{Draft v0.3 --- \today}

\begin{document}
\maketitle

\begin{abstract}
Data selection for language-model pretraining is driven by scalars that
blend every capability together: loss under a reference model, similarity to
a target distribution, or redundancy. We propose scoring data against a
specific \emph{mechanism}: reverse-engineer the circuit a model must build,
reimplement its computation as a cheap program, and grade each document by
how much it exercises that program. We instantiate the principle for
induction heads, whose operation---locate the previous occurrence of the
current context and copy what followed---is the operation of the LZ77
compressor. On synthetic induction tasks, per-document compression gain
predicts the training step at which induction heads form
(Spearman $\rho=-0.98$, three seeds, two architectures), tracks the realised
copy length of the data at $\rho=0.995$ where the generating parameter
tracks it at $0.60$, and selects a training subset that matches a
copy-length oracle exactly. Mechanism tracing shows why: the copy-length
distribution selects among \emph{variants} of the induction circuit that
differ in attention lag, each with a fixed accuracy ceiling, so plateaus
previously read as partial learning are completed circuits of a
data-selected variant. A second compressor quantity, normalised compression
distance between documents, dissociates from gain and predicts
generalisation rather than formation. On real Python code from The Stack,
selecting a training subset by gain under a diversity constraint makes
induction heads form 1.5--2$\times$ sooner in attention and 2--3$\times$
sooner in behaviour than a random subset (three seeds, no overlap), at a
stable cost of 0.2 nats of validation loss that doubling the budget does not
remove; once formed, both circuits copy rare identifiers equally well in
distribution. The compressor's own match-length histogram predicts, before
training, on which corpora selection can act. The complete real-data study
cost under ten dollars of compute.
\end{abstract}

% ======================================================================
\section{Introduction}
% ======================================================================

The data that trains a language model is chosen by heuristics that
optimise a single quantity. Perplexity filters keep documents a reference
model finds neither too easy nor too hard~\citep{wenzek2020ccnet};
importance resampling matches a target
distribution~\citep{xie2023dsir}; domain reweighting learns a mixture from
a proxy model~\citep{xie2023doremi}; learnability-based selection keeps
points that reduce held-out loss~\citep{mindermann2022rho}; deduplication
removes near-copies~\citep{lee2022dedup,abbas2023semdedup,tirumala2023d4}.
These methods treat the model as a black box. They report averages over
benchmarks, and they cannot say which capability a given decision buys or
sells, because they never look at what the model builds.

Mechanistic interpretability has meanwhile produced detailed accounts of
individual circuits---induction heads~\citep{olsson2022induction},
indirect-object identification~\citep{wang2022ioi}, modular
arithmetic~\citep{nanda2023grokking}---and has shown that specific
distributional properties of training data switch some of them on and
off~\citep{chan2022data,reddy2023mechanistic}. That work varies the data
\emph{generator} and observes the model. It does not provide a way to grade
an arbitrary corpus before training.

This paper connects the two. The proposal is simple: \textbf{for a circuit
the model should build, find the cheapest algorithm that performs the same
computation, run it on each document, and use its output as a data score.}
We call this \emph{mechanism-matched scoring}, and we instantiate it once,
completely, for induction heads. The induction head's operation---find the
earlier occurrence of the current token, attend there, copy what
followed---is the operation of the LZ77 compression
algorithm~\citep{ziv1977lz77}, whose match statistics are exposed by any
\texttt{gzip} implementation. The use of compressors to estimate entropy and
cross-entropy of strings was developed
by~\citet{benedetto2002language,baronchelli2005zippers}; we use their two
quantities---per-document compression gain and cross-document compression
distance---and show that they measure two distinct properties of training
data with two distinct consequences for the trained model.

\paragraph{Contributions.}
\begin{enumerate}
\item \textbf{Compression gain predicts induction-head formation}
  (\S\ref{sec:gain}). Across repeat fractions, three seeds and two
  architectures, the formation step is a monotone function of gain
  ($\rho=-0.98$ attention-only, $-0.92$ with MLPs; $-1.0$ on condition
  means).
\item \textbf{The compressor measures the causal variable better than the
  generator} (\S\ref{sec:oracle}). Gain tracks realised copy length at
  $\rho=0.995$; the repeat-fraction parameter that generated the data
  tracks it at $0.60$. Selection by gain matches a copy-length oracle
  (97\% subset overlap, identical formation) and beats selection by the
  generating parameter in three of three pools.
\item \textbf{Copy-length distribution selects the circuit variant}
  (\S\ref{sec:lag}). With MLPs present the induction circuit that forms
  keys on the token $k$ positions back; $k$ is set by the copy-length
  distribution, agrees across seeds, and fixes an exact accuracy ceiling.
  Plateaus at 0.66--0.76 are completed lag-$k$ circuits.
\item \textbf{Compression distance predicts generalisation, not formation}
  (\S\ref{sec:diversity}). The two quantities dissociate (pooled
  $r=-0.27$); with gain fixed, distance predicts held-out generalisation
  ($+0.57$, $+0.75$ partial); with distance fixed, gain predicts neither
  outcome.
\item \textbf{A learned scorer repairs the compressor's blind spot}
  (\S\ref{sec:tiers}). Where accidental matches swamp gain (small
  vocabularies), a reference model with an induction head restores
  prediction ($\rho$ from $-0.72$ to $-0.93$ pooled).
\item \textbf{On real code, selection accelerates formation at a measured
  cost} (\S\ref{sec:stack}). On 2.5B tokens of Python, a gain-plus-diversity
  subset forms induction heads 1.5--2$\times$ sooner in attention and
  2--3$\times$ sooner in behaviour than a random subset, pays 0.2 nats of
  validation loss stably across budgets, and yields an equally load-bearing
  copy head in distribution. The match-length histogram predicts in advance
  where selection can act: on children's stories it cannot, and does not.
\end{enumerate}

A methodological thread runs through the paper. Every synthetic task was
mechanism-traced \emph{before} it was swept, and this was not optional: the
first task we designed to elicit copying was solved by exclusion
(\S\ref{sec:twoname}), and the induction task was redesigned three times as
tracing exposed positional and fixed-offset shortcuts
(Appendix~\ref{app:redesign}). A data-selection claim that skips this step
is a claim about a mechanism the model may not have built.

% ======================================================================
\section{Related Work}
% ======================================================================

\paragraph{Data selection.}
Model-based scores---perplexity~\citep{wenzek2020ccnet}, learnability
\citep{mindermann2022rho}, gradient norms~\citep{paul2021el2n}---and
distribution-based scores---DSIR~\citep{xie2023dsir},
DoReMi~\citep{xie2023doremi}---optimise a global quantity.
Deduplication~\citep{lee2022dedup} and its semantic
variants~\citep{abbas2023semdedup,tirumala2023d4} remove redundancy and
report gains in average loss. Skill-It~\citep{chen2023skillit} orders data
by contribution to evaluation-defined skills. We differ in defining the
target as a circuit and the scorer as a reimplementation of that circuit's
computation, and in reporting circuit-level consequences that averages
hide.

\paragraph{Compression and language models.}
\citet{deletang2023compression} show that language-model loss is a
compression rate under arithmetic coding; \citet{huang2024compression}
report that compression rate correlates with downstream performance as a
global signal; \citet{jiang2023gzip} classify text with gzip distance. We
use compression as a circuit-specific detector and identify the internal
statistic---LZ77 match length---that carries the signal.

\paragraph{Induction heads and training data.}
\citet{olsson2022induction} identify induction heads and their abrupt
formation; \citet{chan2022data} show that burstiness, Zipfian skew and
class count in the data gate in-context learning; \citet{reddy2023mechanistic}
give a mechanistic model of the dynamics. Those works vary the generator
and observe the model. We measure the documents.

% ======================================================================
\section{Method}\label{sec:method}
% ======================================================================

\subsection{Compressor quantities}
Following \citet{baronchelli2005zippers}, the LZ77 code rate of a string
estimates its entropy, and compressing $B$ after $A$ estimates the
cross-entropy $C(A|B)\simeq(L_{A+B}-L_A)/L_B$, where $L_X$ is the
compressed length of $X$. We derive two quantities.

\emph{Gain.} Per-document compressibility beyond unigram statistics,
$\gain(B)=1-L_B/L_{\mathrm{shuf}(B)}$, where $\mathrm{shuf}(B)$ is a token
permutation preserving frequencies. Gain is large when a document contains
long verbatim repeats and zero for a document with no structure beyond its
token frequencies. \todo{confirm normalisation in \texttt{analysis/zipper.py}.}

\emph{Distance.} The normalised compression distance
$\ncd(A,B)=(L_{AB}-\min(L_A,L_B))/\max(L_A,L_B)$~\citep{li2004ncd}.
Corpus diversity is the mean $\ncd$ over a random sample of document pairs
at fixed length.

\emph{Match-length histogram.} LZ77 emits a sequence of (offset, length)
matches. The distribution of match lengths is the mechanistically
meaningful statistic underlying gain, and we use it as a pre-training
diagnostic of whether a corpus contains structure the score can act on.

All quantities are computed on token-id sequences with a sliding-window
LZ77 (zlib), reusing one compressor state for the context when computing
conditional scores.

\subsection{Why LZ77 matches induction}
An induction head~\citep{olsson2022induction} composes a previous-token
head, which writes ``preceded by $x$'' at each position, with a head that at
the current token $x$ attends to positions so labelled and copies the token
found there. LZ77 at each position searches backward for the longest string
that begins here and has already occurred, and emits a pointer to it. Both
locate the previous occurrence of the current context and reuse its
continuation. The compressor's match length is the length of the repeat
the induction head can exploit.

\subsection{Selection}
Given a pool of documents and a token budget, the \emph{gain} selector
ranks by gain and takes the top of the ranking. The \emph{gain + diversity}
selector proceeds greedily, taking the highest-gain candidate whose $\ncd$
to every document already selected exceeds a threshold, which rejects
near-duplicates. \todo{threshold value.} We compare against uniform random
selection and, on synthetic data, against oracles that see the generating
parameters.

% ======================================================================
\section{Synthetic Experiments}\label{sec:synth}
% ======================================================================

\subsection{Tasks}
\paragraph{Induction.} Random token sequences of length $T$ over vocabulary
$V$ containing a repeat of random offset and random length up to
$\ell_{\max}$; the target inside the repeat is the token that followed the
earlier occurrence. Knobs: repeat fraction, $V$, noise, and the copy-length
distribution. Appendix~\ref{app:redesign} records two earlier versions that
tracing showed to be solvable without content matching.
\todo{$T$, $V$, $\ell_{\max}$ defaults.}

\paragraph{Two-name.} Template sentences with two names, one duplicated;
the target is the non-duplicated name, after~\citet{wang2022ioi}. Knobs:
pool size (4--64), Zipf skew, repetition rate, noise. Held-out splits: names
never used as the answer (held-out IO), name pairs never co-occurring
(held-out pairs), and a leak fraction in which held-out names appear as
answers.

\subsection{Models and measurement}
Two-layer four-head transformers with rotary positions by default;
attention-only and with-MLP variants; 3L8H and 4L8H replications. No
weight decay unless stated. Three seeds per condition; checkpoints every
100 steps.

\emph{Formation step.} We use three definitions and report their
agreement: (i) behavioural, the first sustained crossing of an accuracy
threshold; (ii) attention-based, the first checkpoint at which the
lag-aware prefix-matching pattern appears; (iii) patching-based, the first
checkpoint at which patching the circuit's heads recovers $\geq50\%$ of the
clean-minus-corrupt logit difference. Over 39 attention-only runs the
three agree at $\rho=0.99$ with a mean gap of one checkpoint interval.

\emph{Circuit analysis.} Per checkpoint and per head: attention mass to
the relevant positions and direct logit attribution, decomposed exactly
into heads, MLPs, embeddings and the LayerNorm-plus-unembed term.
Faithfulness is accuracy with everything outside the circuit ablated.

\subsection{Gain predicts formation}\label{sec:gain}

\begin{figure}[t]\centering
\todo{\texttt{results/zipper\_vs\_formation.png}, two panels}
\caption{Formation step of the induction head against per-document gain.
Attention-only (left): $\rho=-0.98$ over 12 formed runs. With MLPs (right):
$\rho=-0.92$ over 18 formed runs. $\rho=-1.0$ on condition means in both.
Formation is a jump from accuracy $<0.2$ to $>0.8$ within one or two
100-step intervals; seed spread 0--200 steps.}
\label{fig:gain}
\end{figure}

Table~\ref{tab:gain} lists the repeat-fraction sweep. Below a gain of about
0.07 (attention-only) or 0.10 (MLPs) the circuit does not form within
budget. With MLPs, formation is two to four times later throughout and only
the densest condition snaps cleanly; the others plateau
(\S\ref{sec:lag}). Weight decay at $10^{-3}$ and $10^{-2}$ leaves
formation, lag and ceiling unchanged. A 3L8H attention-only replication
gives $\rho=-0.94$ with the jumps intact.

\begin{table}[t]\centering\small
\begin{tabular}{lrrr}\toprule
Repeat fraction & Gain & Formation, attn-only & Formation, MLP \\\midrule
0.98 & 0.201 & 800   & 1267 \\
0.90 & \todo{} & --- & \todo{} \\
0.85 & \todo{} & --- & \todo{} \\
0.75 & 0.154 & 900   & 4367 \\
0.60 & \todo{} & --- & \todo{} \\
0.50 & 0.101 & 1733  & 6867 \\
0.35 & 0.068 & 3600  & never ($>12$k) \\
0.20 & 0.034 & never ($>6$k) & never \\
0.10 & 0.012 & never & never \\\bottomrule
\end{tabular}
\caption{Formation step (mean of three seeds) by condition. Attention-only
and MLP sweeps used different fraction grids; \todo{fill rows.}}
\label{tab:gain}
\end{table}

\subsection{The compressor measures copy length; the generator does not}
\label{sec:oracle}

In Table~\ref{tab:gain} gain is a deterministic function of the generating
parameter, so \S\ref{sec:gain} shows only that the compressor orders
repetition correctly. A selection experiment separates the two. From a
heterogeneous pool of documents with hidden per-document repeat fractions
we select a fixed token budget by (a) gain, (b) the hidden generating
parameter, (c) realised copy length measured after generation, and (d)
uniformly at random, and train three seeds on each.

Across three independent pools, gain-selected models form the circuit at
step 700; parameter-selected at 800--1000; random at 5200--6000 or never.
The copy-length-selected subset overlaps the gain-selected one by 97\% and
forms at step 700 in all three seeds with final accuracy 0.94--0.95, the
compressor's numbers exactly. The generating parameter is the worse
selector because it describes what was requested rather than produced: it
tracks realised copy length at $\rho=0.60$, admitting documents with high
nominal repetition but short copies (24\% short-copy documents in its
subset, none in the compressor's). Gain tracks realised copy length at
$\rho=0.995$. \todo{Figure: four-arm formation curves.}

\subsection{Copy-length distribution selects the circuit variant}
\label{sec:lag}

\begin{figure}[t]\centering
\todo{lag figure: lag and ceiling against copy-length distribution}
\caption{With MLPs the induction circuit keys on the token $k$ positions
back (lag $k$). Holding $\ell_{\max}=24$: all-long copies give $k=0$,
formation at 800, accuracy 0.997; uniform lengths give $k=19$ and a 0.68
ceiling; mostly-short copies never form. All three seeds agree on $k$ in
every condition.}
\label{fig:lag}
\end{figure}

A lag-$k$ induction circuit keys on the token $k$ positions behind the
current one and copies $k$ positions ahead of the earlier occurrence. It
fires only once the model is at least $k$ tokens into a repeat, so on a
copy-length distribution $p(\ell)$ its accuracy ceiling is
$\mathbb{E}[\max(\ell-k,0)]/\mathbb{E}[\ell]$. \todo{verify against
measured ceilings.}

Every MLP model in our sweeps built a lagged variant, with the lag agreeing
across seeds in every condition. The lag follows the copy-length
distribution (Figure~\ref{fig:lag}), not its maximum and not coverage:
across the repeat-fraction sweep the lag was 12, 15, 19, 21 and 22 for
$\ell_{\max}$ of 16, 19, 24, 27 and 28, but at fixed $\ell_{\max}=24$ the
lag ranges from 0 to 19 depending only on the distribution. Attention-only
models choose lag 0 at every moderate repeat fraction; MLP models never
do. Weight decay does not alter the lag.

Two consequences follow. The plateaus at 0.66--0.76 are not partial
circuits but complete lag-$k$ circuits whose ceiling is predicted by $k$; a
behaviour-only study would misreport them. And the variable that selects
the variant is the variable the compressor measures, which is why
\S\ref{sec:oracle} came out as it did. Why mixed copy lengths favour a
long lag is not yet traced; a boundary-precision account, in which a
lag-$k$ head avoids confident errors at the ends of short copies, predicts
that a geometric length distribution should restore lag 0
(\S\ref{sec:discussion}).

\subsection{Compression distance predicts generalisation, not formation}
\label{sec:diversity}

\begin{figure}[t]\centering
\todo{\texttt{results/diversity\_pooled.png}}
\caption{Partial correlations of gain and $\ncd$ with each two-name outcome
over 45 runs, each quantity controlling for the other.}
\label{fig:div}
\end{figure}

On the two-name task, repetition rate raises gain (0.10 to 0.34) with
$\ncd$ flat (0.67--0.73), and Zipf skew lowers $\ncd$ (0.728 to 0.647) with
gain flat (0.215--0.236); pooled with the pool-size sweep the two correlate
at $r=-0.27$. Table~\ref{tab:div} gives partial correlations.

\begin{table}[t]\centering\small
\begin{tabular}{lrr}\toprule
Outcome & $\ncd$ (partial) & Gain (partial) \\\midrule
Formation, attention-based & $-0.04$ & $+0.03$ \\
Formation, patching-based  & $+0.16$ & $\phantom{+}0.00$ \\
Held-out pairs accuracy    & $+0.57$ & $-0.20$ \\
Held-out IO accuracy       & $+0.29$ & $+0.09$ \\
Crossing of held-out logit difference & $+0.75$ & $+0.21$ \\\bottomrule
\end{tabular}
\caption{Diversity predicts generalisation; gain adds nothing once
diversity is fixed; neither predicts formation of the exclusion circuit.}
\label{tab:div}
\end{table}

Formation is flat across every knob except the extremes. At repetition
rate 1.0 every later sentence repeats an earlier one, an induction shortcut
pre-empts the task circuit, and formation is delayed from 333 to 1500
steps: the regime in which data admits a cheaper mechanism and the model
builds that instead of the target.

\subsection{The two-name circuit is exclusion, not copying}
\label{sec:twoname}

The task was designed to elicit name-copying after~\citet{wang2022ioi}. The
model built something else. Two layer-1 heads attend to the duplicated
name with mass 0.9--1.0 and push its logit down; no head places more than
0.2 attention on the answer. The answer's margin of 10--12 logits over
pool names absent from the prompt comes from MLP~1 and MLP~0 fed by
layer-0 attention, not from a name-mover. Patching the two heads recovers
86--100\% of the logit difference; ablating them collapses held-out-IO
accuracy to zero in every run and validation accuracy to 0.5 in five of
twelve, a third head serving as backup in the rest. In a 2L4H model the
circuit is most of the model (6--8 of 8 heads for threshold recovery).

\paragraph{Priors and their cost.}
Names never used as answers acquire a prior against being the answer,
carried by MLP~0 (the unembed term is $<0.1$ everywhere). The prior forms
\emph{during} circuit formation---the held-out logit difference deepens to
about $-4.5$ over the first 250 leaked examples---then climbs linearly at
about 0.01 logits per leaked example and crosses zero after 600--1000
examples, roughly 1\% of sentences seen. Zipf skew weakens the prior faster
than it costs pair generalisation (held-out IO 0.36 to 0.63; crossing 2600
to 267 steps).

\subsection{A learned scorer repairs the compressor's blind spot}
\label{sec:tiers}

At vocabulary 16, accidental short matches compress the shuffled baseline
as well as the document, and gain collapses (0.021 versus 0.101 at
vocabulary 100) while formation does not. A fixed reference model with an
induction head scores the two datasets alike (0.090 versus 0.104). Pooled
over 39 induction runs, Spearman with formation is $-0.72$ for gain and
$-0.93$ for the reference model; the entire gap is the vocabulary family,
and within every other family the two agree. Gain and the model's loss are
the same quantity under two compressors~\citep{deletang2023compression};
the learned one wins exactly where LZ77's matching is uninformative.
\todo{\texttt{results/tiers\_vs\_formation.png}.}

% ======================================================================
\section{Real Corpora}\label{sec:real}
% ======================================================================

\subsection{Children's stories: the diagnostic predicts a null}
\label{sec:tinystories}

On a 65M-token slice of TinyStories no induction head forms in any of seven
runs; the prefix-matching score never leaves its initialisation baseline.
The match-length histogram says why: median LZ77 match of 2 tokens,
0.1--0.35\% of tokens in matches of $\geq8$, against 27\% in the synthetic
subset that forms at step 700. Selection moves bigram repeats from 12\% to
18\% and long matches by tenths of a percent. Natural-prose induction rides
on single-token repeats (47\% of tokens), which the shuffled-baseline gain
discounts, and reported prose phase changes sit at $10^9$--$10^{10}$
tokens~\citep{olsson2022induction}, 15--150 times this budget. The score
has no range here and the model has no budget here; the histogram
identified the first before any model was trained.

\subsection{Python code: selection accelerates formation}
\label{sec:stack}

\paragraph{Setup.}
A 2.5B-token Python slice of The Stack~\citep{kocetkov2022stack}, cut into
windows of 128 tokens. \todo{tokenizer; window count.} The match-length
histogram of random windows already carries more long-match structure than
the synthetic subset of \S\ref{sec:oracle}, and the gain-ranked top fifth
more than double it. Arms: gain-plus-diversity selection and uniform random
selection of matched size, a 4-layer model, 491M training tokens (30,000
steps), three seeds each, checkpoints every 500--1000 steps; one seed per
arm extended to 983M tokens. Measures: the prefix-matching score
of~\citet{olsson2022induction} and induction accuracy on repeated random
sequences the model has not seen; validation loss on held-out code; and an
in-context delta on rare identifiers (below). \todo{model width, batch,
LR.}

A preliminary local experiment on 64M tokens of installed-package Python,
repeated over epochs, grew a head five times stronger in the selected arm
but formed in neither; fresh data at the same budget, not more steps, was
the fix, as with TinyStories (Appendix~\ref{app:local}).

\paragraph{Formation.}
\begin{table}[t]\centering\small
\begin{tabular}{llrrrrr}\toprule
Arm & Seed & $\pm\geq0.5$ & Ind.\ acc $\geq0.5$ & Final $\pm$ & Final ind.\ acc & Val loss \\\midrule
Gain + div. & 0 & 1,000 & 2,000 & 0.75 & 0.86 & 1.80 \\
Gain + div. & 1 & 1,000 & 2,000 & 0.89 & 0.86 & --- \\
Gain + div. & 2 & 1,000 & 2,000 & 0.83 & 0.84 & 1.78 \\
Random      & 0 & 2,000 & 6,000 & 0.82 & 0.70 & 1.60 \\
Random      & 1 & 1,500 & 4,000 & 0.74 & 0.78 & --- \\
Random      & 2 & 2,000 & 5,000 & 0.74 & 0.73 & 1.59 \\\midrule
Gain + div., 983M & 0 & 1,000 & 2,000 & 0.75 & 0.86 & 1.745 \\
Random, 983M      & 0 & 2,000 & 5,500 & 0.84 & 0.70 & 1.557 \\\bottomrule
\end{tabular}
\caption{The Stack, 491M tokens unless stated. Formation steps are the
first sustained crossing of 0.5. \todo{fill missing val-loss cells.}}
\label{tab:stack}
\end{table}

\begin{figure}[t]\centering
\todo{\texttt{results/stack/stack\_trajectories.png}}
\caption{Prefix-matching score and induction accuracy by step for both
arms, three seeds, with validation loss. The selected arm snaps
(prefix-matching 0.02 to 0.71 within the first 1,000 steps; induction
accuracy 0.14 to 0.60 between 1,000 and 2,000); the random arm seeps
(prefix-matching 0.31 to 0.70 over steps 1,000--4,000; induction accuracy
0.17 to 0.64 over 2,000--8,000).}
\label{fig:stack}
\end{figure}

The selected arm forms the head at step 1,000 in every seed and the
behaviour at step 2,000 in every seed; the random arm at 1,500--2,000 and
4,000--6,000 (Table~\ref{tab:stack}): 1.5--2$\times$ sooner in attention and
2--3$\times$ sooner in behaviour, with no overlap on any measure. The
selected arm's induction accuracy stays 0.06--0.16 above the random arm's
to the end of training while prefix-matching overlaps (0.75--0.89 versus
0.74--0.84): both heads attend where they should, and the selected model
executes the copy more accurately on random sequences. Doubling the budget
leaves formation and final induction accuracy unchanged in both arms and
improves validation loss by 0.05 in both; the 0.2-nat gap between arms
persists. The trajectories reproduce, between two selections of one real
corpus, the snap-versus-seep contrast observed between the induction and
two-name tasks (Figure~\ref{fig:stack}).

\paragraph{In-context delta.}
Absolute loss on repeated identifiers orders like validation loss and so
measures distribution quality, not the head. We instead measure, for each
rare identifier (not among the 2,000 most frequent names) occurring at
least twice in a validation window, the cross-entropy at its first
occurrence $\mathrm{CE}_1$ and at its second $\mathrm{CE}_2$; the delta
$\mathrm{CE}_1-\mathrm{CE}_2$ is the in-context benefit, and distribution
quality cancels. As controls we compute the same delta for repeated common
tokens and a name-swap test in which the identifier is replaced throughout
the window's prefix by a substitute.

\begin{table}[t]\centering\small
\begin{tabular}{llrrrrr}\toprule
Arm & Seed & $\mathrm{CE}_1$ & $\mathrm{CE}_2$ & Delta (frac.\ $>0$) & Common-token delta & Prefers swap \\\midrule
Gain + div. & 0 & 9.14 & 4.56 & 4.58 (0.86) & 0.74 & 0.57 \\
Gain + div. & 1 & 9.06 & 4.78 & 4.28 (0.84) & 0.73 & 0.57 \\
Gain + div. & 2 & 9.07 & 4.54 & 4.53 (0.86) & 0.73 & 0.58 \\
Random      & 0 & 8.39 & 3.98 & 4.41 (0.88) & 0.66 & 0.54 \\
Random      & 1 & 8.20 & 4.02 & 4.18 (0.87) & 0.66 & 0.47 \\
Random      & 2 & 8.29 & 3.73 & 4.56 (0.88) & 0.66 & 0.54 \\\midrule
Gain + div., 983M & 0 & 9.20 & 4.27 & 4.93 (0.86) & 0.72 & 0.60 \\
Random, 983M      & 0 & 8.16 & 3.59 & 4.56 (0.88) & 0.66 & 0.55 \\\bottomrule
\end{tabular}
\caption{In-context delta on 4,807 rare-identifier pairs and 20,867
repeated-common-token pairs from 2,000 validation windows, final
checkpoints. Delta declines with occurrence distance in every arm alike
(5.1--5.5 nats at $\leq32$ tokens; 2.7--3.0 at 128--256).}
\label{tab:delta}
\end{table}

The copy head is load-bearing in distribution in both arms: 4.2--4.9 nats
on rare identifiers, positive for 84--88\% of them, six times the delta on
common tokens, and the swap control moves 5--6 nats off the original name
and 7.4--8.9 onto the substitute. At the end of training the two arms copy
about equally (mean delta 4.46 versus 4.38 over three seeds, within seed
spread); $\mathrm{CE}_1$ orders like validation loss, the distribution cost
again. At 983M tokens the selected arm's delta is 4.93 versus 4.56, a
modest edge on one seed.

\paragraph{Reading.}
On real code, selection by the compressor buys \emph{when} the induction
circuit forms and a circuit that copies more accurately on
syntax-free sequences; it does not buy in-distribution copying once both
circuits exist, and it costs 0.2 nats on the distribution that doubling
the budget does not recover. The synthetic finding transfers---gain
governs formation and a gain-heavy selection pays on the distribution---and
the recommendation is unchanged: use both compressor quantities.

\paragraph{Trade-off curve at small scale.}
On the local 64M-token pool, four arms at matched size---gain only, gain
plus diversity, half gain and half random, random---order monotonically on
validation loss (2.43, 2.30, 2.13, 2.03) while the diversity-constrained
arm grows the strongest head of any (prefix-matching 0.10--0.12 versus
0.03--0.07 for gain only). Rejecting near-duplicates removes
cross-document repetition, which an in-context copy head cannot use, and
keeps within-document repetition, which it can. The trade-off is tunable,
and the two quantities set its two axes. \todo{repeat the four arms on The
Stack if budget allows.}

% ======================================================================
\section{Discussion}\label{sec:discussion}
% ======================================================================

\paragraph{The principle.}
Gzip predicts induction not because compression measures data quality in
general but because LZ77 and the induction head perform the same
computation. The generalisation is not ``use gzip'' but ``for each circuit,
find the cheapest program that computes what the circuit computes.'' Where
no such program is at hand, a small model that has already learned the
circuit is a learned implementation of it (\S\ref{sec:tiers}).

\paragraph{What averages hide.}
Deduplication lowers gain; our results predict, and \S\ref{sec:stack}
shows the converse of, a delay in induction. Outlier filtering lowers
diversity; \S\ref{sec:diversity} predicts a cost in rare-item
generalisation and stronger never-an-answer priors. Neither cost appears
in a loss average. Mechanism-matched scores make them measurable before
training, and the match-length histogram says on which corpora a given
score can act at all.

\paragraph{Data picks the variant.}
That the copy-length distribution sets a circuit parameter with a built-in
ceiling (\S\ref{sec:lag}), and that two selections of one real corpus
yield circuits that copy equally in distribution but unequally on
syntax-free sequences (\S\ref{sec:stack}), suggest a broader question:
which data statistics set which circuit parameters. We have one law and
one real-data echo of it.

\paragraph{Limitations.}
Models of 2--4 layers; one circuit family; one learning rate. The
real-data acceleration is measured at formation steps of one to six
thousand, early relative to any production budget; the persistent
difference is in the circuit's generality, not its in-distribution use.
The lag mechanism is described but not explained. The exclusion circuit's
formation is insensitive to every data property varied, so the two-name
task tests diversity, not formation. Prose corpora at our budget are below
the induction phase change, and our diagnostic correctly reports that the
score has no range there, but we have not shown selection acting on prose.

\paragraph{Future work.}
A second circuit where the hand-written scorer should fail and the learned
one should win: synthetic facts with paraphrases, where gain selects
against the data that builds retrieval and $\ncd$ selects for it. Tracing
the lag mechanism through boundary loss and a geometric length
distribution. The four-arm trade-off on The Stack. Interference between
circuits on mixed data.

\paragraph{Reproducibility.}
All synthetic experiments run on a laptop CPU or its integrated GPU. The
complete real-code study---data download, tokenisation, scoring, eight
training runs to 491M--983M tokens, and analysis---cost about \$6.60 of
rented GPU time. \todo{repo link.}

\bibliographystyle{plainnat}
\bibliography{refs}

\appendix
\section{Induction-task redesigns}\label{app:redesign}
Three versions were built; each was mechanism-traced before sweeping.
(1) Fixed repeat offset: solved by a layer-0 positional head with no content
matching. (2) Random offset, fixed length: content matching present, but
keyed on ``the token $r$ positions back'' because a fixed copy length
privileges that feature. (3) Random offset and length: no shortcut
remains; used throughout. Learned absolute positions, not weight decay,
were found to block formation in an early configuration.

\section{Local code pool}\label{app:local}
12,037 installed-package Python files, 498k windows of 128 tokens; top
100k windows by gain versus three random 100k subsets; 2L4H, 8,000 and
24,000 steps. Gain of subset 0.365 versus 0.162; 45\% versus 14\% of tokens
in matches $\geq8$. Prefix-matching at 8,000 steps 0.03--0.07 (rising from
step 5,000) versus 0.01 flat; at 24,000 steps 0.17 versus 0.03; validation
loss 2.41 versus 1.90. Four-arm results in \S\ref{sec:stack}.

\section{Additional figures}
\todo{\texttt{leak\_sample\_efficiency.png}; per-run patching summaries;
tier comparison.}

\end{document}
